%=============================================================================
% CONFIGURACIÓN DE FORMATO BAJO NORMAS APA 7ma EDICIÓN (RIGUROSO)
%=============================================================================

%-----------------------------------------------------------------------------
% DISEÑO DE PÁGINA Y MÁRGENES (APA 7: 1 pulgada o 2.54 cm en los 4 lados)
%-----------------------------------------------------------------------------
\usepackage[a4paper, margin=2.54cm]{geometry}

%-----------------------------------------------------------------------------
% NUMERACIÓN DE PÁGINAS Y ENCABEZADOS (Solución al problema de la esquina superior)
%-----------------------------------------------------------------------------
\usepackage{fancyhdr}
\pagestyle{fancy}
\fancyhf{} % Borra por completo encabezados y pies de página anteriores

% Configura el número de página estrictamente en la esquina superior derecha
\rhead{\thepage} 
\lhead{}  % En trabajos estudiantiles NO se usa encabezado de texto (Running Head)
\cfoot{}  % Garantiza que no quede ningún número flotando abajo al centro

% Elimina las líneas horizontales decorativas del encabezado y pie (Prohibidas en APA)
\renewcommand{\headrulewidth}{0pt}
\renewcommand{\footrulewidth}{0pt}

% REDEFINICIÓN CRÍTICA DEL ESTILO 'PLAIN'
% Evita que LaTeX mueva el número abajo en el índice o inicios de sección
\fancypagestyle{plain}{
	\fancyhf{}
	\rhead{\thepage}
	\renewcommand{\headrulewidth}{0pt}
	\renewcommand{\footrulewidth}{0pt}
}

%-----------------------------------------------------------------------------
% INTERLINEADO Y SANGRÍAS
%-----------------------------------------------------------------------------
\usepackage{setspace}
\doublespacing % APA 7 exige doble espacio reglamentario en todo el texto, tablas y notas
\setlength{\parindent}{1.27cm} % Sangría reglamentaria de 0.5 pulgadas (1.27 cm)

%-----------------------------------------------------------------------------
% IDIOMA Y TRADUCCIÓN
%-----------------------------------------------------------------------------
\usepackage[spanish, es-tabla]{babel} % 'es-tabla' evita que LaTeX escriba "Cuadro" en vez de "Tabla"

%-----------------------------------------------------------------------------
% FUENTES Y MATEMÁTICAS
%-----------------------------------------------------------------------------
\usepackage[T1]{fontenc}
\usepackage[utf8]{inputenc}
\usepackage{amsmath, amsfonts, amssymb, bm, mathtools}

%-----------------------------------------------------------------------------
% JERARQUÍA DE TÍTULO / ENCABEZADOS (NIVELES APA 7)
%-----------------------------------------------------------------------------
\usepackage{titlesec}

% Nota institucional: APA puro no numera los títulos. Como en ingeniería de la UNT 
% se exige numeración decimal (1., 1.1), se conserva el formato numérico pero adaptado 
% visualmente a la posición y estilo estricto de APA 7.

% Nivel 1 (\section): Centrado, Negrita
\titleformat{\section}[block]
{\normalfont\large\bfseries\centering}{\thesection.}{1em}{}

% Nivel 2 (\subsection): Alineado a la izquierda, Negrita
\titleformat{\subsection}[block]
{\normalfont\normalsize\bfseries}{\thesubsection.}{1em}{}

% Nivel 3 (\subsubsection): Alineado a la izquierda, Negrita, Cursiva
\titleformat{\subsubsection}[block]
{\normalfont\normalsize\bfseries\itshape}{\thesubsubsection.}{1em}{}

% Nivel 4 (\paragraph): Con sangría, Negrita, punto al final, texto continuo en la misma línea
\titleformat{\paragraph}[runin]
{\normalfont\normalsize\bfseries}{\theparagraph.}{1em}{}[.]
\titlespacing*{\paragraph}{\parindent}{1.5ex plus 0.5ex minus 0.2ex}{1em}

% Nivel 5 (\subparagraph): Con sangría, Negrita, Cursiva, punto al final, texto continuo en la misma línea
\titleformat{\subparagraph}[runin]
{\normalfont\normalsize\bfseries\itshape}{\thesubparagraph.}{1em}{}[.]
\titlespacing*{\subparagraph}{\parindent}{1.5ex plus 0.5ex minus 0.2ex}{1em}

%-----------------------------------------------------------------------------
% TABLAS Y FIGURAS AUTOMATIZADAS (ESTILO APA 7)
%-----------------------------------------------------------------------------
\usepackage{graphicx}
\usepackage{longtable}
\usepackage{array}
\usepackage{booktabs} % Indispensable para tablas APA (elimina líneas verticales)
\usepackage{float}
\usepackage{caption}

% Automatización completa de los Captions:
% El número ("Tabla 1") saldrá en NEGRITA. El título descriptivo irá ABAJO en una nueva línea y en CURSIVA.
\captionsetup{
	labelsep=newline,          % Envía el título a una línea nueva
	justification=raggedright, % Alineado a la izquierda
	singlelinecheck=false,     % Evita que captions cortos se centren solos
	labelfont=bf,              % Etiqueta en Negrita
	textfont=it                % Texto del título en Cursiva
}

%-----------------------------------------------------------------------------
% OTROS PAQUETES DE TU PROYECTO
%-----------------------------------------------------------------------------
\usepackage{pdfpages}
\usepackage{tocloft}
\usepackage{tikz}
\usepackage[dvipsnames]{xcolor}

%-----------------------------------------------------------------------------
% HIPERVÍNCULOS
%-----------------------------------------------------------------------------
\usepackage[breaklinks]{hyperref}
\hypersetup{
	colorlinks=true,
	linkcolor=black,
	filecolor=black,      
	urlcolor=blue,
	citecolor=black
}